\documentclass[12pt,a4paper]{article}
\usepackage[utf8]{inputenc}
\usepackage[spanish]{babel}
\usepackage{amsmath}
\usepackage{array}
\usepackage{geometry}
\usepackage{booktabs}   
\usepackage{float}      
\usepackage{xltabular}  

\geometry{margin=1in}


\newcommand{\PreserveBackslash}[1]{\let\temp=\\#1\let\\=\temp}
\newcolumntype{C}[1]{>{\PreserveBackslash\centering}p{#1}}
\newcolumntype{L}[1]{>{\PreserveBackslash\raggedright}p{#1}}
\newcolumntype{R}[1]{>{\PreserveBackslash\raggedleft}p{#1}}


\renewcommand{\arraystretch}{1.3}

\title{Resumen de Arquitectura MIPS}
\author{Eric Guacache}
\date{}

\begin{document}
	
	\maketitle
	
	\section*{Nombres de registros, usos y convenio de llamadas}
	
	\begin{table}[H]
		\centering
		\small
		\begin{tabularx}{\textwidth}{l | c | X | c}
			\toprule
			\textbf{NOMBRE} & \textbf{NÚMERO} & \textbf{USO} & \textbf{¿SE CONSERVA?} \\
			\midrule
			\texttt{\$zero} & 0 & Valor constante 0 & N/A \\
			\midrule
			\texttt{\$at} & 1 & Ensamblador temporal & No \\
			\midrule
			\texttt{\$v0-\$v1} & 2-3 & Valores de resultado de funciones y evaluación de expresiones & No \\
			\midrule
			\texttt{\$a0-\$a3} & 4-7 & Argumentos & No \\
			\midrule
			\texttt{\$t0-\$t7} & 8-15 & Temporales & No \\
			\midrule
			\texttt{\$s0-\$s7} & 16-23 & Temporales almacenados & Sí \\
			\midrule
			\texttt{\$t8-\$t9} & 24-25 & Temporales & No \\
			\midrule
			\texttt{\$k0-\$k1} & 26-27 & Reservados para el núcleo del Sistema Operativo & No \\
			\midrule
			\texttt{\$gp} & 28 & Puntero global & Sí \\
			\midrule
			\texttt{\$sp} & 29 & Puntero de pila & Sí \\
			\midrule
			\texttt{\$fp} & 30 & Puntero de marco & Sí \\
			\midrule
			\texttt{\$ra} & 31 & Dirección de retorno & Sí \\
			\bottomrule
		\end{tabularx}
	\end{table}
	
	\section*{Formatos básicos de instrucción}
	
	\begin{table}[H] 
		\centering
		\begin{tabular}{c|c|c|c|c|c|c}
			\toprule
			\textbf{Tipo} & \multicolumn{5}{c}{\textbf{Campos}} & \\
			\midrule
			R & cod oper (31-26) & rs (25-21) & rt (20-16) & rd (15-11) & desplaz (10-6) & func (5-0) \\
			\midrule
			I & cod oper (31-26) & rs (25-21) & rt (20-16) & \multicolumn{3}{c}{inmediato (15-0)} \\
			\midrule
			J & cod oper (31-26) & \multicolumn{5}{c}{dirección (25-0)} \\
			\bottomrule
		\end{tabular}
	\end{table}
	
	\clearpage
	
	\section*{Funciones de lenguaje ensamblador MIPS}
	

	\small
	\begin{xltabular}{\textwidth}{>{\raggedright\arraybackslash}p{2.2cm} | >{\raggedright\arraybackslash}p{2.2cm} | >{\raggedright\arraybackslash}p{3cm} | >{\raggedright\arraybackslash}p{3.5cm} X}
		\caption{Instrucciones MIPS.}\\
		\toprule
		\textbf{Categoría} & \textbf{Instrucción} & \textbf{Ejemplo} & \textbf{Significado} & \textbf{Comentarios} \\
		\midrule
		\endfirsthead 
		
		\multicolumn{5}{c}{{\bfseries \tablename\ \thetable{} -- Continuación de la página anterior}} \\
		\toprule
		\textbf{Categoría} & \textbf{Instrucción} & \textbf{Ejemplo} & \textbf{Significado} & \textbf{Comentarios} \\
		\midrule
		\endhead 
		
		\bottomrule
		\endfoot 
		\bottomrule
		\endlastfoot 
		
		Aritmética 
		& add & \texttt{add \$s1, \$s2, \$s3} & \texttt{\$s1 = \$s2 + \$s3} & Tres operandos; datos en registros \\ \cline{2-5}
		& subtract & \texttt{sub \$s1, \$s2, \$s3} & \texttt{\$s1 = \$s2 - \$s3} & Tres operandos; datos en registros \\ \cline{2-5}
		& add immed. & \texttt{addi \$s1, \$s2, 100} & \texttt{\$s1 = \$s2 + 100} & Usado para sumar constantes \\
		\midrule 
		Transferencia de dato 
		& load word & \texttt{lw \$s1, 100(\$s2)} & \texttt{\$s1 = Mem[\$s2+100]} & Palabra de memoria a registro \\ \cline{2-5}
		& store word & \texttt{sw \$s1, 100(\$s2)} & \texttt{Mem[\$s2+100] = \$s1} & Palabra de registro a memoria \\ \cline{2-5}
		& load half & \texttt{lh \$s1, 100(\$s2)} & \texttt{\$s1 = Mem[\$s2+100]} & Media palabra de memoria a registro \\ \cline{2-5}
		& store half & \texttt{sh \$s1, 100(\$s2)} & \texttt{Mem[\$s2+100] = \$s1} & Media palabra de registro a memoria \\ \cline{2-5}
		& load byte & \texttt{lb \$s1, 100(\$s2)} & \texttt{\$s1 = Mem[\$s2+100]} & Byte de memoria a registro \\ \cline{2-5}
		& store byte & \texttt{sb \$s1, 100(\$s2)} & \texttt{Mem[\$s2+100] = \$s1} & Byte de registro a memoria \\ \cline{2-5}
		& load upper immed. & \texttt{lui \$s1, 100} & \texttt{\$s1 = 100 $\cdot 2^{16}$} & Cargar constante en los 16 bits de mayor peso \\
		\midrule 
		Lógica 
		& and & \texttt{and \$s1, \$s2, \$s3} & \texttt{\$s1 = \$s2 \& \$s3} & Tres registros operandos; AND bit-a-bit \\ \cline{2-5}
		& or & \texttt{or \$s1, \$s2, \$s3} & \texttt{\$s1 = \$s2 | \$s3} & Tres registros operandos; OR bit-a-bit \\ \cline{2-5}
		& nor & \texttt{nor \$s1, \$s2, \$s3} & \texttt{\$s1 = $\sim$(\$s2 | \$s3)} & Tres registros operandos; NOR bit-a-bit \\ \cline{2-5}
		& and immed. & \texttt{andi \$s1, \$s2, 100} & \texttt{\$s1 = \$s2 \& 100} & AND bit-a-bit registro con constante \\ \cline{2-5}
		& or immed. & \texttt{ori \$s1, \$s2, 100} & \texttt{\$s1 = \$s2 | 100} & OR bit-a-bit registro con constante \\ \cline{2-5}
		& shift left & \texttt{sll \$s1, \$s2, 10} & \texttt{\$s1 = \$s2 $\ll$ 10} & Desplazamiento a la izquierda por constante \\ \cline{2-5}
		& shift right & \texttt{srl \$s1, \$s2, 10} & \texttt{\$s1 = \$s2 $\gg$ 10} & Desplazamiento a la derecha por constante \\
		\midrule 
		Salto condicional 
		& branch on eq & \texttt{beq \$s1, \$s2, L} & if (\texttt{\$s1 == \$s2}) go L & Comprueba igualdad y bifurca relativo al PC \\ \cline{2-5}
		& branch not eq & \texttt{bne \$s1, \$s2, L} & if (\texttt{\$s1 != \$s2}) go L & Comprueba si no igual y bifurca relativo al PC \\ \cline{2-5}
		& set on less & \texttt{slt \$s1, \$s2, \$s3} & if (\texttt{\$s2 < \$s3}) \texttt{\$s1=1} else \texttt{\$s1=0} & Compara si es menor que; usado con beq, bne \\ \cline{2-5}
		& set on less imm & \texttt{slti \$s1, \$s2, 100}& if (\texttt{\$s2 < 100}) \texttt{\$s1=1} else \texttt{\$s1=0} & Compara si es menor que una constante \\
		\midrule 
		Salto incond. 
		& jump & \texttt{j 2500} & go to 10000 & Salto a la dirección destino \\ \cline{2-5}
		& jump reg & \texttt{jr \$ra} & go to \texttt{\$ra} & Para retorno de procedimiento \\ \cline{2-5}
		& jump and link & \texttt{jal 2500} & \texttt{\$ra = PC+4}; go 10000 & Para llamada a procedimiento \\
	\end{xltabular}
	
\end{document}

